\chapter{Appendix IV: Methodological Protocol and Reproducibility Checklist}
\label{app:methodological-protocol}

\section{Objective of This Appendix}
This appendix provides a replayable protocol for independent verification. The goal is to reduce ambiguity in experimental setup, data handling, and evaluation workflow.

\section{Environment Profile}
\begin{table}[ht]
\centering
\caption{Execution Environment Matrix}
{\small
\begin{tabular}{@{}L{2.8cm}L{3.8cm}L{3.2cm}@{}}
\toprule
\textbf{Layer} & \textbf{Reference Environment} & \textbf{Purpose} \\ \midrule
Mobile app build & Android Gradle toolchain + Kotlin/JVM & Edge pipeline compilation and test execution \\
Analytics prototype & Python environment under project-local setup & Dataset exploration and algorithm calibration \\
Analytics API runtime & Docker multi-stage image & Service-level consistency and deployment parity \\
Thesis build & LaTeX engine (containerized pass recommended) & Deterministic manuscript generation \\
\bottomrule
\end{tabular}
}
\end{table}

\section{Data Handling Protocol}
\subsection{Acquisition and Normalization}
\begin{enumerate}
    \item Import public datasets into structured local folders.
    \item Validate schema-level consistency (image labels, coordinate formats, units).
    \item Normalize timestamps, coordinate precision, and confidence ranges.
    \item Remove near-duplicate image samples and overlapping temporal windows.
\end{enumerate}

\subsection{Split Strategy}
\begin{enumerate}
    \item Use route-level grouping to avoid leakage between training and validation.
    \item Stratify by lighting and speed context where available.
    \item Keep a held-out subset for final reporting only.
\end{enumerate}

\section{Calibration Workflow}
\subsection{Sensor Calibration}
\begin{enumerate}
    \item Initialize Kalman parameters ($Q$, $R$) from literature-backed defaults.
    \item Run short calibration windows on representative road profiles.
    \item Adjust anomaly thresholds by minimizing false alerts while preserving impact recall.
\end{enumerate}

\subsection{Fusion Calibration}
\begin{enumerate}
    \item Start from baseline weights $(\alpha,\beta,\gamma)$.
    \item Evaluate under day/night and low/high-speed buckets.
    \item Validate that weight normalization and threshold semantics remain stable.
\end{enumerate}

\section{Evaluation Workflow}
\subsection{Core Metrics}
Required metrics for each mode:
\begin{equation}
Precision = \frac{TP}{TP+FP}, \quad Recall = \frac{TP}{TP+FN}, \quad F1 = \frac{2PR}{P+R}.
\end{equation}

\subsection{Comparative Modes}
The comparative evaluation must include:
\begin{enumerate}
    \item CV-only,
    \item Sensor-only,
    \item Fixed Fusion,
    \item Adaptive Fusion.
\end{enumerate}

\subsection{Acceptance Rule}
A fusion update is accepted only if:
\begin{enumerate}
    \item the F1 gain is positive versus CV-only baseline,
    \item no integrity regressions appear in timestamp handling,
    \item no invalid geolocation artifacts are introduced.
\end{enumerate}

\section{Post-Audit Regression Gate}
After each critical change, execute a focused gate:
\begin{itemize}
    \item Fusion unit tests (including timestamp compatibility),
    \item app compilation check,
    \item thesis recompilation to keep manuscript aligned with implementation.
\end{itemize}

\subsection{Web Access and Authentication Preflight}
Before running dashboard rehearsals, execute a dedicated preflight to separate
service-availability faults from credential-configuration faults:
\begin{enumerate}
    \item verify that \path{/login} is reachable on the expected preview endpoint,
    \item verify that Firebase configuration is not using placeholder values,
    \item verify that the production build output target is \path{dist-app} (to avoid legacy permission conflicts on \path{dist}),
    \item when using containerized checks, run the tooling profile before full rehearsal execution.
\end{enumerate}

This sequence prevents false incident attribution (for example, treating a local
service outage as an authentication policy bug) and improves replay consistency.

\subsection{Rollback and Troubleshooting Baseline}
For reproducible incident handling, each rehearsal/deploy cycle should include a
minimal rollback and troubleshooting checklist:
\begin{enumerate}
    \item snapshot current manifests/scripts before applying changes,
    \item if a regression appears, restore the previous known-good revision and rerun the gate set,
    \item classify failures by layer (Android, web, API, RBAC, infrastructure) before mitigation,
    \item record root cause and mitigation in the rehearsal log for future replay.
\end{enumerate}

This discipline reduces recovery time and avoids mixing infrastructure faults,
credential faults, and application regressions in a single undifferentiated incident.

\subsection{Live Infrastructure Closure Checklist}
Before promoting infrastructure portability claims from code-level verification to
environment-level verification, collect a minimal live evidence set:
\begin{enumerate}
    \item apply OpenShift overlay manifests and capture rollout status for both services,
    \item verify Route admission and endpoint reachability (web root and API health),
    \item archive pod readiness snapshots and failure-free execution window,
    \item attach these artifacts to the preflight package used for final closure.
\end{enumerate}

This checklist ensures that portability claims are grounded in runtime behavior,
not only in manifest presence.

\subsection{Final Rollout Gate Sequence (Step 1--Step 3)}
For final closure, apply the following sequence in order:
\begin{enumerate}
    \item \textbf{Step 1 -- OpenShift hardening and live validation}: execute the live checklist above; if \path{oc}/\path{kubectl} are unavailable, classify the step as \emph{environment-blocked} rather than a code regression.
    \item \textbf{Step 2 -- Login reproducibility}: lock one authentication mode (Auth Emulator in this thesis baseline) and archive positive/negative/RBAC traces.
    \item \textbf{Step 3 -- Web bundle reproducibility}: execute production build and archive chunk summary and exit status (including post-fix logs when mitigation is applied).
\end{enumerate}

This stepwise rule prevents mixing infrastructure-access limitations with application-level regressions and keeps closure decisions auditable.

\subsection{Closure Execution Note (2026-04-22)}
The final closure run followed the Step 1--Step 3 sequence exactly and archived an additional evidence set under \path{output/statistics/final_*}. During this pass:
\begin{itemize}
    \item Step 2 and Step 3 were re-validated successfully (auth/RBAC and web build reproducibility).
    \item API smoke and Android regression gates were re-executed and stored in dedicated final logs.
    \item Step 1 remained explicitly tagged as \emph{environment-blocked} because OpenShift CLIs were unavailable on the host.
\end{itemize}

This distinction preserves methodological rigor: software-functional closure is evidenced independently from infrastructure-access constraints.

\section{Artifact Checklist for External Examiner}
\begin{table}[ht]
\centering
\caption{Minimal Reproducibility Artifacts}
{\small
\begin{tabular}{@{}L{4.0cm}L{6.2cm}@{}}
\toprule
\textbf{Artifact} & \textbf{Verification Question} \\ \midrule
Calibration constants table & Are numerical defaults explicitly declared? \\
Fusion equations and thresholds & Is decision logic mathematically traceable? \\
Comparative evaluation tables & Are baseline and adaptive modes both reported? \\
Regression test evidence & Are critical fixes protected by automated checks? \\
Post-fix web build log & Does the build complete successfully after mitigation? \\
Deployment descriptors & Can runtime topology be reconstructed? \\
\bottomrule
\end{tabular}
}
\end{table}

\section{Replication Notes}
A faithful replication should preserve:
\begin{enumerate}
    \item sensor sampling assumptions,
    \item confidence threshold semantics,
    \item context-aware weighting logic,
    \item route-level split discipline,
    \item human-in-the-loop validation policy.
\end{enumerate}

If these constraints are changed, results remain valid only as a new experimental variant and must be reported as such.
